% Generated draft: requires mathematical and source review.
\documentclass[11pt]{article}
\usepackage[margin=25mm]{geometry}
\usepackage{fontspec}
\setmainfont{Latin Modern Roman}
\usepackage{amsmath,amssymb,mathtools,unicode-math}
\setmathfont{Latin Modern Math}
\usepackage{enumitem,needspace,xurl,hyperref}
\setlength{\parindent}{0pt}
\setlength{\parskip}{5pt}
\setlength{\emergencystretch}{3em}
\newcommand{\N}{\mathbb{N}}
\newcommand{\Z}{\mathbb{Z}}
\newcommand{\Q}{\mathbb{Q}}
\newcommand{\R}{\mathbb{R}}
\newcommand{\C}{\mathbb{C}}
\newcommand{\F}{\mathbb{F}}
\newcommand{\A}{\mathbb{A}}
\newcommand{\PP}{\mathbb{P}}
\newcommand{\E}{\mathbb{E}}
\newcommand{\eps}{\varepsilon}
\newcommand{\dd}{\,\mathrm{d}}
\begin{document}
\section*{UnsolvedMath ID: 30005566}
\textbf{Source problem code:} catalogue entry
\subsection{Definitions and mathematical statement}
Let \(Y\) be a smooth projective variety and let \(D\subset Y\) be a simple-normal-crossings divisor. The source does not specify the ground field or characteristic; write \(k\) for the unspecified ground field, without imposing or asserting characteristic-independent validity of any intended construction. Under the usual reduced-divisor convention, write \(D=\bigcup_{i\in I}D_i\), where the \(D_i\) are its irreducible components. Simple normal crossings means that, geometrically and locally in the étale topology, \(D\) is a union of distinct coordinate hyperplanes. For \(J\subseteq I\), put
\[
D_J=\bigcap_{j\in J}D_j,\qquad D_{\varnothing}=Y,\qquad D_J^{\circ}=D_J\setminus\bigcup_{i\notin J}D_i.
\]
The nonempty connected components of the locally closed sets \(D_J^{\circ}\) will be called the strata of the pair.

The ordinary Hilbert functor of \(Y\) assigns to a \(k\)-scheme \(T\) the closed subschemes \(Z\subseteq Y\times_k T\) that are flat and of finite presentation over \(T\). It is represented by the ordinary Hilbert scheme \(\operatorname{Hilb}(Y)\). A fixed-Hilbert-polynomial version additionally requires a choice of polarization and polynomial; neither is specified in the source. The ordinary Hilbert scheme depends on \(Y\), so specifying \(D\) supplies additional incidence or logarithmic data to be studied on an enhancement of this moduli problem.

For context, the pair has the divisorial logarithmic structure
\[
j:Y\setminus D\hookrightarrow Y,\qquad
M_Y=\mathcal O_Y\mathbin{\times}_{j_*\mathcal O_{Y\setminus D}}j_*\mathcal O_{Y\setminus D}^{\times},
\]
with its natural map to \(\mathcal O_Y\), interpreted on the étale site. This defines the logarithmic structure associated with the boundary, but does not by itself specify a logarithmic Hilbert moduli functor.

\textbf{Source problem.} For the pair \((Y,D)\), construct and characterize a logarithmic Hilbert moduli object encoding how subschemes meet \(D\) and its strata, and explain its relationship with \(\operatorname{Hilb}(Y)\). This retains both the reconstructed request for a logarithmic Hilbert scheme and the original question, ``How does the presence of D affect the Hilbert scheme Hilb(Y )?''

\textbf{Unresolved specifications.} The supplied record does not define the objects and morphisms of the intended logarithmic moduli problem, its logarithmic test objects, admissibility or transversality conditions, contact or tropical data, numerical invariants, or whether modifications or expansions of the target are allowed. It also does not specify whether the resulting object should be a scheme, algebraic space, stack, or logarithmic analogue, or what precise comparison with the ordinary Hilbert scheme is requested. Thus the source gives a construction-and-characterization question, rather than a uniquely determined representability theorem or conjecture. The definitions above provide notation and do not resolve these missing specifications.
\subsection{Short English statement}
For a smooth projective variety \(Y\) with a simple-normal-crossings boundary \(D\), seek a logarithmic Hilbert moduli object that records how subschemes meet the boundary strata and relates this information to \(\operatorname{Hilb}(Y)\). The intended moduli functor and precise characterization remain unspecified. The supplied record labels the problem partially solved; that dated assessment is not independently verified here.
\subsection{Sources}
\begin{itemize}
\item Supplied local UnsolvedMath registry record, permanent ID \(30005566\), title \emph{Logarithmic Hilbert Schemes for Divisor-Stratified Varieties}, source problem code \texttt{catalogue entry}, in the collection \emph{Oberwolfach Reports Open Problems}; proposed year recorded as 2023. No separate proposer is recorded.
\item Supplied workshop provenance: \emph{Recent Trends in Algebraic Geometry} (2023), Oberwolfach Reports, DOI \texttt{10.4171/owr/2023/27}, pp.~1485--1542. \url{https://doi.org/10.4171/owr/2023/27}. The supplied background describes these reports as documenting mathematical workshops and their open-problem discussions.
\item Supplied original wording includes \emph{The Hilbert scheme in logarithmic geometry, after Kennedy-Hunt Dhruv Ranganathan}. This wording is retained as attribution evidence; its punctuation and the respective attribution roles are not resolved by the record.
\item Kennedy-Hunt, \emph{The logarithmic Quot space: foundations and tropicalisation}. Supplied literature link: \url{https://arxiv.org/abs/2308.14470}.
\item Kennedy-Hunt--Ranganathan, \emph{Logarithmic Quot spaces, boundedness, and K-tropicalizations}. Supplied literature link: \url{https://arxiv.org/abs/2508.08175}.
\item Literature assessment dated exactly \texttt{2026-08-22}: the supplied curation labels the problem \texttt{partially\_solved} and reports that Kennedy-Hunt constructs the logarithmic Hilbert/Quot moduli problem and proves foundational separatedness, properness, boundedness, and tropicalization results, while the broad effect of \(D\) on the classical Hilbert scheme is not a single fully settled classification. This is supplied dated curation derived from the OWR material, not an independently verified literature assessment; its status and cited claims require independent verification.
\item The registry labels the statement \texttt{corrected\_verified} and supplies the verification note that it was reconstructed and checked against the official Oberwolfach report. These are supplied provenance claims, not verification performed for this draft. The registry records its last update as \texttt{2026-09-16T12:06:53Z}.
\end{itemize}
\end{document}
